\section{The empty divisorial locus}
\label{fol:section}

The obstruction considered in this section is the absence of a meromorphic
pluricanonical section, including one with poles.  We keep this distinction
throughout: a rational line bundle on a nonalgebraic compact complex space
need not have a global meromorphic frame.

\begin{proposition}\label{fol:empty-case}
Assume Assumption~\ref{hyp:campana}.  Let $Y$ be a normal compact K\"ahler
fourfold with klt singularities, and suppose that the actual rational
canonical line bundle $K_Y$ is analytically nef.  Suppose that $Y$ contains
no prime divisors.  Let $p\colon M\to Y$ be a smooth compact K\"ahler log
resolution.  If $a(M)=0$ and $K_M$ is analytically pseudo-effective, then
some positive tensor power of $K_M$ has a nonzero meromorphic section.
\end{proposition}

We prove the proposition by contradiction.  Until the end of the section,
assume its hypotheses and, in addition, assume that
\begin{equation}\label{fol:no-meromorphic}
 H^0\bigl(M,\mathcal M_M\otimes\OO_M(mK_M)\bigr)=0
 \qquad\text{for every integer }m>0,
\end{equation}
where $\mathcal M_M$ is the sheaf of meromorphic functions.  Put
\begin{equation}\label{fol:canonical-comparison}
 L=p^*K_Y,\qquad K_M=L+C,\qquad \alpha=2\pi c_1(L).
\end{equation}
Here $C$ is the rational exceptional discrepancy divisor, and the middle
identity is an identity of actual rational line bundles, using the
canonical comparison on the isomorphism locus of $p$.  No sign is imposed
on $C$.  Every prime divisor on $M$ is $p$-exceptional, because $Y$ has
none.  We use the weight convention $\ddc=\ii\partial\bar\partial$, so a
metric weight on $L$ has curvature class $\alpha$.

\subsection{Virtual vanishings and holomorphic forms}

The starting Euler-characteristic/Hodge-theoretic route to holomorphic
forms was informed by the related forms strategy of Vikash \cite{Vikash2026}.
That work considers smooth minimal fourfolds without effective divisors
or surfaces.  The present normal klt setting has different hypotheses,
and the required vanishings and forms construction are proved below.

\begin{lemma}\label{fol:virtual-vanishing}
Let $f\colon N\to M$ be a proper surjective generically finite morphism,
where $N$ is a connected smooth compact K\"ahler manifold.  Then
\begin{equation}\label{fol:virtual-data}
 a(N)=q(N)=0,\qquad K_N\text{ is pseudo-effective}.
\end{equation}
and no positive tensor power of $K_N$ has a nonzero meromorphic section.
These conclusions also hold after any further smooth compact K\"ahler
modification or finite cover followed by resolution.
\end{lemma}

\begin{proof}
First consider meromorphic functions and tensors.  A proper generically
finite map of normal complex spaces factors, by Stein factorization, as a
proper modification followed by a finite map.  Meromorphic functions on a
modification descend to the normal target, and finite maps have local
meromorphic traces and norms.  These constructions use local meromorphic
function fields and therefore do not require a global meromorphic frame.

If $u$ is a meromorphic function on $N$, its characteristic polynomial over
$M$, computed on the finite part of the factorization, has meromorphic
coefficients on $M$.  These coefficients are constant because $a(M)=0$.
Thus $u$ satisfies a polynomial with constant coefficients.  Connectedness
then makes $u$ constant, and $a(N)=0$.

Suppose that $s$ is a nonzero meromorphic section of $mK_N$.  Over a
coordinate neighborhood $U\subset M$, choose a nowhere vanishing local
frame $\tau$ of $mK_M$.  The differential pullback $f^*\tau$ is a
holomorphic pluricanonical tensor, generically nonzero, so
$s/(f^*\tau)$ is a meromorphic function over $f^{-1}U$; zeros of the
Jacobian merely contribute a meromorphic divisor to this quotient.
If $d$ is the generic degree of $f$, its norm is the coefficient of a
meromorphic section of $mdK_M$: replacing $\tau$ by $g\tau$ divides the
norm by $g^d$.  The local sections consequently glue and are nonzero.
This contradicts \eqref{fol:no-meromorphic}.  The ordinary Jacobian
formula
\[
 K_N=f^*K_M+R_f,\qquad R_f\geq0,
\]
proves pseudo-effectivity of $K_N$.

It remains to prove $q(N)=0$.  If $q(N)>0$, the Albanese map has a
positive-dimensional image.  Let $h\colon N\to B$ be its Stein
factorization onto its normal image.  The space $B$ is finite over a
subvariety of a torus.  On a smooth compact K\"ahler model $B'$ of $B$,
the pullbacks of suitable ambient one-forms have a nonzero wedge of
degree $\dim B$.  Hence $\kappa(B')\geq0$.  On every neat smooth model
used to compute the invariant base in Assumption~\ref{hyp:campana}, the
orbifold boundary is effective; birational invariance of ordinary
Kodaira dimension therefore gives
\[
 \kappa(h\mid0)\geq\kappa(B')\geq0.
\]
This comparison concerns the invariant base term, not a boundary on an
arbitrarily chosen initial model.

Choose a positive singular metric on $K_N$.  Its local potentials restrict
to almost every smooth fiber of $h$, by local integrability and Fubini's
theorem.  We may choose such a fiber outside the countably many proper
analytic subsets excluded by the very-general-fiber convention in
Assumption~\ref{hyp:campana}.  If $F$ is that fiber, then $K_F$ is
pseudo-effective and $\dim F\leq3$.  Ordinary canonical nonvanishing for
smooth compact K\"ahler manifolds of dimension at most three gives
$\kappa(F)\geq0$.  In dimension three this follows from the K\"ahler
threefold minimal model theorem and abundance for nef normal compact
K\"ahler threefold pairs; see
\cite{HPMinimalModels,DasOuAbundanceII}.  In dimensions at most two it is
the classical curve and surface statement.  Applying
Assumption~\ref{hyp:campana} to $(N,0)$ and $h$ gives
\[
 \kappa(N)\geq\kappa(F)+\kappa(h\mid0)\geq0,
\]
contrary to the absence of a meromorphic pluricanonical section on $N$.
The same proof applies to any further morphism of the stated kind.
\end{proof}

In particular, $q(M)=0$.  Moreover,
\begin{equation}\label{fol:nonzero-alpha}
 \alpha\ne0.
\end{equation}
Indeed, if $c_1(L)=0$ in real cohomology, a multiple of $L$ has torsion
integral first Chern class.  After a further multiple that class vanishes.
The exponential sequence and $H^1(M,\OO_M)=0$ make that line bundle
trivial.  Equation~\eqref{fol:canonical-comparison} would then give a
meromorphic section of a positive multiple of $K_M$.

\begin{lemma}\label{fol:eventual-vanishing}
For every holomorphic vector bundle $\mathcal V$ on $M$,
\[
 H^0(M,\mathcal V\otimes\OO_M(mL))=0
\]
for every sufficiently large positive integer $m$ divisible by the index
of $L$.  Furthermore,
\begin{equation}\label{fol:forms-dimensions}
 \chi(M,\OO_M)=0,\qquad h^{2,0}(M)\geq1,\qquad h^{3,0}(M)\geq2.
\end{equation}
\end{lemma}

\begin{proof}
Choose an integer $r>0$ such that $A=rL$ is a line bundle.  A section of
$\mathcal V\otimes A^k$ is a morphism $A^{-k}\to\mathcal V$.  Among
all nonzero sections with $k>0$, choose a tuple
$s_1,\ldots,s_t$, of twists $k_1,\ldots,k_t$, which is generically
independent and has maximal possible cardinality.  If there are no such
sections the assertion is immediate.  Otherwise $t\leq\rk\mathcal V$.
Let $\mathcal W\subset\mathcal V$ be the saturation of the image of the
corresponding direct sum of line bundles.  Every further section lies in
$\mathcal W$ generically, by maximality, and hence everywhere as a map
into this saturated subsheaf.

A further nonzero section $s$ of twist $k$ replaces at least one $s_i$
while preserving generic independence.  The two determinants are nonzero
sections, respectively, of
\[
 \det\mathcal W\otimes A^{\sum k_j}
 \quad\text{and}\quad
 \det\mathcal W\otimes A^{\sum k_j+k-k_i},
\]
where determinants are reflexive.  Their ratio is a nonzero meromorphic
section of $A^{k-k_i}$.  For $k>\max_i k_i$, this is a positive multiple
of $L$.  After clearing denominators in $C$, multiplication by its
canonical meromorphic divisor section converts it into a meromorphic
section of a positive multiple of $K_M$, contradicting
\eqref{fol:no-meromorphic}.  This proves the eventual vanishing without
choosing a meromorphic frame of $A$.

By Lemma~\ref{met:zero-lelong}, $L$ has a semipositive metric with zero
Lelong numbers everywhere.  The multiplier ideal of every positive
integral multiple of this metric is trivial by Skoda integrability
\cite{SkodaIntegrability}.  The hard Lefschetz theorem with multiplier
ideals \cite[Theorem~0.1]{DPSHardLefschetz} therefore gives, for every
$0\leq j\leq4$, a surjection
\[
 H^0(M,\Omega_M^{4-j}\otimes\OO_M(mL))
 \longrightarrow H^j(M,\OO_M(K_M+mL))
\]
for divisible $m>0$.  The source is zero for all sufficiently large such
$m$.  Riemann--Roch makes $\chi(M,\OO_M(K_M+mL))$ a polynomial in $m$;
its vanishing on all these multiples implies that its constant term
$\chi(M,\OO_M(K_M))$ is zero.  Serre duality in dimension four yields
$\chi(M,\OO_M)=0$.

Equation~\eqref{fol:no-meromorphic} gives $h^{4,0}=0$, and
Lemma~\ref{fol:virtual-vanishing} gives $h^{1,0}=0$.  If $h^{2,0}=0$,
then all of $H^2(M,\R)$ is of type $(1,1)$.  Approximating a K\"ahler
class by a rational one and applying the Kodaira embedding theorem would
make $M$ projective, contrary to $a(M)=0$.  Thus $h^{2,0}\geq1$.
Finally, Hodge symmetry and the Euler characteristic identity give
$h^{3,0}=1+h^{2,0}\geq2$.
\end{proof}

\subsection{A saturated integrable conormal line}

Write $K=K_M$.  Choose $0\ne\sigma\in H^0(M,\Omega_M^2)$.  Holomorphic
forms on $M$ are closed, and $\sigma^2=0$ because it is a canonical
section.  Thus $\sigma$ has rank two wherever it is nonzero.  The kernel
of $\sigma$ is integrable there, and its saturated conormal is a rank-two
subsheaf
\[
 F\subset\Omega_M^1,\qquad G=\Omega_M^1/F.
\]
All determinants and line factors below are taken reflexively.  The form
$\sigma$ is a nonzero section of $\det F$, so
\begin{equation}\label{fol:F-and-G}
 \det F=\OO_M(H),\qquad \det G=\OO_M(K-H)
\end{equation}
for an effective divisor $H$.

We use the natural identification
$\Omega_M^3\simeq K\otimes T_M$, obtained by contraction with a local
volume form.  Global linearly independent three-forms are generically
independent: coefficients of a dependence relative to a maximal
generically independent tuple are global meromorphic functions, hence
constants.  The same observation applies to their projections into
$F^*\otimes K$.

If $v$ and $w$ are the $K$-valued vector fields associated with two
three-forms, then $\sigma(v,w)$ is a section of $2K$, so it is zero.
Since the form induced by $\sigma$ on the rank-two quotient of $T_M$ is
nondegenerate at a general point, the projections of all these vector
fields to $F^*\otimes K$ have rank at most one.  On the other hand, two
generically independent tangent vector fields would have a nonzero
wedge in
\[
 2K\otimes(\det G)^*=\OO_M(K+H).
\]
Removing the divisor factor would give a prohibited meromorphic
canonical section.  The space of global three-forms has dimension at
least two, so we can choose associated vector fields $v_1,v_2$ such that
$v_1$ is nonzero and tangent to the kernel of $\sigma$, while $v_2$ is
not tangent.  Denote their three-forms by $\beta_1,\beta_2$.

Let $E\subset F$ be the saturation of the line defined by the nonzero
$K$-valued one-form
\[
 \boldsymbol\eta=\iota_{v_2}\sigma.
\]
It follows that
\begin{equation}\label{fol:E-class}
 E=\OO_M(-K+e)
\end{equation}
for a divisor $e$ (its divisorial zero divisor).  Saturation makes $E$ a
reflexive rank-one sheaf, hence a line bundle on $M$.  Its inclusion in
$\Omega_M^1$ is a subbundle away from a set of codimension at least two.
Likewise, all the rank-two sheaves and quotients just used are vector
bundles at general points of every prime divisor.

The rank-two identity
\[
 K\otimes G^*\simeq\OO_M(H)\otimes G
\]
turns $v_1$ into a section of $\OO_M(H)\otimes G$.  Saturating its image
produces a line factor $\OO_M(j-H)\subset G$, where $j$ is a divisor;
by the determinant identity the other factor is $\OO_M(K-j)$.

\begin{lemma}\label{fol:Bott-integrability}
The line $E$ is an integrable saturated conormal.  On a dense open set,
its line is the intersection of the degree-one wedge annihilators of
$\sigma$ and $\beta_2$.
\end{lemma}

\begin{proof}
Work first on the complement of a codimension-at-least-two set where
$F,G,E$ and the line factors above are bundles.  Let $D\subset T_M$ be
the integrable distribution annihilated by $F$.  Lie derivative along
$D$ defines the partial Bott connection on $F$.  Its second fundamental
form for $E$ is a morphism
\begin{equation}\label{fol:Bott-tensor}
 B_E\colon E\longrightarrow(F/E)\otimes D^*
                 =(F/E)\otimes G.
\end{equation}
For clarity, this is tensorial in both variables: for $w\in D$ and a
local section $s$ of $E$, the terms introduced by replacing $s$ by $fs$
are multiples of $s$, while the additional term in $\mathcal L_{fw}s$
is $s(w)\,\dd f=0$.

If the composite of $B_E$ with the second line factor of $G$ is nonzero,
it is a nonzero section of the line bundle of class
\[
 (H-E)+(K-j)-E=3K+H-j-2e.
\]
If that composite vanishes and $B_E\ne0$, it factors through the first
line factor and gives a nonzero section of class
\[
 (H-E)+(j-H)-E=2K+j-2e.
\]
These sections extend across the omitted set by reflexivity and Hartogs'
theorem.  Removing their divisor factors gives a nonzero meromorphic
section of $3K$ or $2K$, respectively, contradicting
\eqref{fol:no-meromorphic}.  Thus $B_E=0$.

In local Frobenius coordinates for $D$, write its transverse coordinates
as $(z_1,z_2)$ and a frame of $E$ as $a\,\dd z_1+b\,\dd z_2$.
Vanishing of $B_E$ says that the ratio $a/b$, where defined, is constant
along the plaques of $D$.  After multiplication by an invertible local
function, the frame is therefore a one-form on the two-dimensional
transversal.  Every rank-one conormal on a smooth surface is integrable.
The resulting Frobenius identity extends to the entire regular locus
of $E$ by holomorphic equality, including general points of divisorial
zeros of the original forms.

The wedge annihilator of $\sigma$ in degree one is $F$.  The annihilator
of $\beta_2$ is the hyperplane of one-forms vanishing on $v_2$.  Since
$v_2$ is not tangent to $D$, their intersection is a line.  It contains
$\iota_{v_2}\sigma$, which both wedges to zero with $\sigma$ and
evaluates to zero on $v_2$.  Hence that line is $E$.
\end{proof}

\begin{lemma}\label{fol:annihilation}
For every closed positive $(1,1)$-current $T$ in the class $\alpha$, one has
\begin{equation}\label{fol:annihilations}
 T\wedge\sigma\wedge\bar\sigma=0,\qquad
 T\wedge\beta_2\wedge\bar\beta_2=0
\end{equation}
on the open subset where $p$ is an isomorphism onto a smooth open subset
of $Y$.  In particular,
\begin{equation}\label{fol:T-eta}
 T\wedge\boldsymbol\eta=0
\end{equation}
on a dense open subset with analytic complement.
\end{lemma}

\begin{proof}
Off the singular set of the saturated codimension-one foliation, whose
codimension is at least two, holomorphic first integrals give frames
$\dd z$ for $E$.  If $\dd z_a=g_{ab}\dd z_b$ on an overlap, then
$\dd\log g_{ab}$ belongs to $E$: differentiating the transition identity
gives $\dd g_{ab}\wedge\dd z_b=0$.  Consequently the first Chern cocycle
wedges to zero with both $\sigma$ and $\beta_2$ there.

These vanishings extend in Dolbeault cohomology to all of $M$.  Indeed,
for a locally free sheaf $\mathcal V$ on a smooth manifold and an analytic
set $A$ of codimension at least two,
$\mathcal H^j_A(\mathcal V)=0$ for $j=0,1$, by depth and Hartogs'
theorem.  The local-to-global sequence for cohomology with support
therefore gives an injection
$H^1(M,\mathcal V)\to H^1(M\setminus A,\mathcal V)$.
Apply it to $\Omega_M^3$ and $\Omega_M^4$.  We obtain
\begin{equation}\label{fol:Bott-Chern}
 c_1(E)\wedge[\sigma]=0,\qquad
 c_1(E)\wedge[\beta_2]=0.
\end{equation}

We next dispose of the exceptional divisor terms.  Fix a K\"ahler form
$\omega_Y$ on $Y$.  For every exceptional prime $P$ on $M$,
\begin{equation}\label{fol:exceptional-integrals}
 \int_P \sigma\wedge\bar\sigma\wedge p^*\omega_Y=0,\qquad
 \int_P \ii\beta_2\wedge\bar\beta_2=0.
\end{equation}
Here and below integration on a prime divisor can be computed on a
smooth resolution.  To see the assertion, rationality of klt
singularities gives $Rp_*\OO_M=\OO_Y$, hence
$H^k(Y,\OO_Y)\simeq H^k(M,\OO_M)$.  Resolve $P$ and the image $p(P)$,
and resolve the graph so that the resulting smooth compact K\"ahler
space $\widehat P$ maps to a smooth model $B$ of $p(P)$.  Functoriality
of restriction makes the classes of $\bar\sigma|_{\widehat P}$ and
$\bar\beta_2|_{\widehat P}$ pullbacks from $H^2(B,\OO_B)$ and
$H^3(B,\OO_B)$, respectively.  Conjugation and compact K\"ahler Hodge
theory give the corresponding pullback assertions for their de Rham
classes.  The class of $p^*\omega_Y$ restricts from $B$ as well.
Since $\dim B\leq2$, the degree-six products in
\eqref{fol:exceptional-integrals} vanish.

By \eqref{fol:E-class} and \eqref{fol:canonical-comparison},
$c_1(E)=-c_1(L)$ modulo exceptional divisor classes.  Combining
\eqref{fol:Bott-Chern} and \eqref{fol:exceptional-integrals} gives
\[
 \int_M T\wedge\sigma\wedge\bar\sigma\wedge p^*\omega_Y=0,
 \qquad
 \int_M T\wedge\ii\beta_2\wedge\bar\beta_2=0.
\]
The integrands are nonnegative: $\sigma$ is decomposable wherever
nonzero, and every three-form on a four-dimensional vector space is
decomposable.  The form $p^*\omega_Y$ is strictly positive on the
isomorphism locus.  Thus the zero-mass equalities imply
\eqref{fol:annihilations} there.

This positivity argument applies to currents as well as smooth forms.
Locally, write the coefficient matrix of $T$ as a positive semidefinite
matrix of Radon--Nikodym derivatives with respect to its trace measure.
The two wedge equalities imply, almost everywhere for that measure, that
each one-form direction of the matrix annihilates both $\sigma$ and
$\beta_2$ by wedge product.  Lemma~\ref{fol:Bott-integrability} identifies
their intersection with $E$.  This proves \eqref{fol:T-eta}.
\end{proof}

\subsection{Index charts and a closed transverse current}

We will repeatedly use the following local integrability fact.  All
integrals in its proof are computed in fixed smooth coordinate volumes.

\begin{lemma}\label{fol:Jacobian-integrability}
Let $f\colon V\to U$ be a proper generically finite holomorphic map
between smooth complex manifolds of the same dimension.  Let $\phi$ be
a plurisubharmonic function on $U$ such that
$e^{-s\phi}\in L^1_{\mathrm{loc}}(U)$ for every $s>0$.  Then
\[
 e^{-s(\phi\circ f)}\in L^1_{\mathrm{loc}}(V)
 \qquad(s>0).
\]
The assertion is locally uniform for a family of functions for which
all the indicated exponential integrals downstairs are locally uniformly
bounded.
\end{lemma}

\begin{proof}
Take relatively compact coordinate charts upstairs and downstairs so that
$f$ maps the former into the latter.  Let $J$ be the local holomorphic
Jacobian determinant; it is not identically zero.  A sufficiently small
negative power of $|J|$ is locally integrable.  Choose $P>1$ large enough
that $|J|^{-2/(P-1)}$ is integrable on the chosen compact set.  H\"older's
inequality gives
\begin{equation}\label{fol:Jacobian-Holder}
 \int e^{-s\phi\circ f}
 \leq
 \left(\int e^{-sP\phi\circ f}|J|^2\right)^{1/P}
 \left(\int |J|^{-2/(P-1)}\right)^{(P-1)/P}.
\end{equation}
Change of variables bounds the first factor by the degree of $f$ times
the corresponding exponential integral downstairs, up to fixed
coordinate-volume constants.  Covering compact sets by finitely many
charts proves both assertions.
\end{proof}

A positive closed $(1,1)$-current with zero Lelong numbers has no mass
on any proper analytic subset.  We recall precisely the form used here.
The Skoda--El Mir extension theorem says that restricting such a current
off an analytic set and then extending by zero gives a closed current.
The difference is a positive closed current supported on that set.  The
support theorem reduces it to its divisorial components, whose
coefficients are the corresponding generic Lelong numbers, so the
difference is zero.  These current-theoretic facts, together with Skoda
integrability and Siu decomposition below, are used in their ordinary
smooth-manifold forms; see \cite{DemaillyComplexAnalyticGeometry}.

\begin{lemma}\label{fol:transverse-current}
Let $T$ be a positive current in $\alpha$ with zero Lelong numbers
everywhere.  Choose local plurisubharmonic metric weights $\phi$ for
$T$ on $L$.  On the isomorphism locus of $p$, write
$\boldsymbol\eta=\eta\otimes\tau$ in the corresponding local canonical
frame $\tau$.  The expression
\begin{equation}\label{fol:Theta-definition}
 \Theta_T=\ii e^{-\phi}\eta\wedge\bar\eta
\end{equation}
has a canonical extension to a nonzero positive closed $(1,1)$-current
on $M$.  This extension has no mass on proper analytic subsets and satisfies
\begin{equation}\label{fol:Theta-orthogonality}
 \alpha\smile[\Theta_T]=0.
\end{equation}
The weights are understood with a fixed global additive normalization
when $\Theta_T$ is compared for different currents.
\end{lemma}

\begin{proof}
We first specify the extension through the singular model.  Cover $Y$
by neighborhoods $U$ on which a Cartier pluricanonical power has a
nowhere vanishing frame.  Take the ordinary local index cover
$\widehat U\to U$ defined by a root of that frame, and normalize.
The cover is quasi-\'etale and klt, by the index-cover discrepancy
formula; in particular it has rational singularities.  On its smooth
locus the root is a canonical frame.

Resolve a main component of
$\widehat U\times_U p^{-1}(U)$, obtaining a diagram
\[
 \begin{tikzcd}
 \widetilde U \arrow[r] \arrow[d,"f"'] & \widehat U \arrow[d]\\
 p^{-1}(U) \arrow[r] & U .
 \end{tikzcd}
\]
Here $f$ is proper and generically finite.  The resolution is chosen
isomorphic over the dense smooth open where the index cover is
unramified and $p$ is an isomorphism.  The coefficient $\eta$ in the
root frame is an ordinary holomorphic one-form on that open.
It extends reflexively to $\widehat U$ by Hartogs' theorem, since the
complement of the relevant smooth locus downstairs has codimension at
least two.  The extension theorem for rational complex spaces then
extends it holomorphically to $\widetilde U$; we use
\cite[Corollary~1.8]{KebekusSchnellExtension}.
We continue to write $\eta$ for this holomorphic form and $\phi$ for the
pulled metric weight in the root frame.

Zero Lelong numbers downstairs give
$e^{-s\phi}\in L^1_{\mathrm{loc}}$ for every $s>0$.  By
Lemma~\ref{fol:Jacobian-integrability}, the same is true on
$\widetilde U$.  In particular the pulled weight has zero Lelong
numbers, and the pulled curvature has no analytic-subset mass.
The upstairs expression in \eqref{fol:Theta-definition} has
$L^r_{\mathrm{loc}}$ coefficients for every finite $r$.  This holds
also after multiplication by $\phi$: on a compact set $\phi$ is
bounded above, while $|\phi|e^{-\phi}$ is bounded by a constant times
$1+e^{-2\phi}$.

Define $\Theta_T$ on $p^{-1}(U)$ as $f_*$ of this expression divided
by $\deg f$.  On the good open this is the original formula.  It is
independent of the root frame: if $\tau'=g\tau$, then
$\eta'=g^{-1}\eta$ and $\phi'=\phi-\log|g|^2$.
A pushforward of the indicated $L^r$ coefficient forms has no mass on
the analytic complement, because its inverse image is a proper
analytic subset and hence has zero smooth volume upstairs.
Consequently these local pushforwards agree on overlaps and define
a global positive current.  It is nonzero since
$\boldsymbol\eta$ is generically nonzero.

We now prove closedness, rather than assuming that a transverse-looking
density is closed.  On $\widetilde U$, the saturation of the line defined
by $\eta$ gives an integrable codimension-one conormal.  Integrability
holds first on the good open by Lemma~\ref{fol:Bott-integrability} and
then everywhere by holomorphic equality.  Its singular set $A$ has
codimension at least two.  On a regular foliated chart, including one
through a divisorial zero of $\eta$, write
\begin{equation}\label{fol:local-foliation}
 \eta=f_0\,\dd z
\end{equation}
with $z$ a holomorphic submersion and $f_0$ holomorphic.
Lemma~\ref{fol:annihilation} gives
$\ddc\phi\wedge\dd z=0$ on the good part where $f_0\ne0$.
It holds on the whole chart because the positive closed curvature
has no mass on the remaining analytic set.

In coordinates $(z,w_1,w_2,w_3)$, positivity shows that the only
coefficient of $\ddc\phi$ is its
$\ii\,\dd z\wedge\dd\bar z$ coefficient.  Closedness of that current
makes this coefficient independent, as a distribution, of every $w_j$
and $\bar w_j$.  Thus it is the pullback of a positive measure on the
$z$-disc.  Taking a one-variable subharmonic potential and then solving
the pluriharmonic difference on a smaller polydisc gives
\begin{equation}\label{fol:split-weight}
 \phi=\phi_0(z)+g+\bar g,\qquad g\text{ holomorphic}.
\end{equation}
With $h=f_0e^{-g}$, the upstairs current becomes
\[
 \ii e^{-\phi_0(z)}|h(z,w)|^2\,\dd z\wedge\dd\bar z.
\]
Its $\ddc$ is positive.  More explicitly, only differentiation in the
$w$ directions survives wedging with
$\dd z\wedge\dd\bar z$, and
\begin{equation}\label{fol:leaf-Hessian}
 \ddc\Theta_T
 =
 e^{-\phi_0(z)}
 \ii\partial_w h\wedge\bar\partial_w\bar h
 \wedge\ii\,\dd z\wedge\dd\bar z
 \ \geq\ 0
\end{equation}
on the regular foliated chart.  This is a distributional identity;
the one-variable coefficient is locally integrable, and no transverse
derivative of it contributes to the displayed wedge.

We give the extension estimate across $A$.  On a relatively compact
coordinate ball, choose smooth cutoffs $\chi_\epsilon$ equal to zero
on an $\epsilon$-tube around $A$ and equal to one outside a
$3\epsilon$-tube, with
\[
 |\nabla\chi_\epsilon|\leq C\epsilon^{-1},
 \qquad
 |\nabla^2\chi_\epsilon|\leq C\epsilon^{-2}.
\]
They may be constructed by smoothing the distance cutoffs.  The tubes
have volume $O(\epsilon^4)$.  To justify the bound also when $A$ is
singular, use locally finite area for its analytic components and
monotonicity to bound the number of disjoint balls of radius
$\epsilon/4$ centered on each component by
$C\epsilon^{-2d}$, where $d\leq2$ is its complex dimension.
Enlarging this cover to the tube gives volume
$O(\epsilon^{8-2d})=O(\epsilon^4)$; a compact subset meets finitely
many local components.

If $B$ denotes either $\Theta_T$ or $\phi\Theta_T$ upstairs and
$r>2$, then for any fixed smooth test form the second-derivative
cutoff terms are bounded by
\begin{equation}\label{fol:cutoff-estimate}
 C\epsilon^{-2}\|B\|_{L^r}
       \operatorname{vol}(N_{3\epsilon}(A))^{1-1/r}
 \leq C_r\|B\|_{L^r}\epsilon^{\,2-4/r}
 \longrightarrow0.
\end{equation}
Terms with only one derivative are smaller.  Testing
\eqref{fol:leaf-Hessian} against $\chi_\epsilon$ times a positive
test form and passing to the limit therefore proves
$\ddc\Theta_T\geq0$ on all of $\widetilde U$.

Proper pushforward commutes with $\ddc$, so $\ddc\Theta_T$ is a
global positive exact $(2,2)$-current on $M$.  Pairing with the square
of a K\"ahler form gives zero; positivity then makes it the zero
current.  On the good open the chart maps are local biholomorphisms.
All branch contributions to the pushforward are positive, so each
upstairs Hessian in \eqref{fol:leaf-Hessian} vanishes there separately.
Since $e^{-\phi_0}>0$ almost everywhere, the holomorphic derivatives
$\partial_w h$ vanish there, and hence on every regular foliated chart
by holomorphic continuation.  Thus
\begin{equation}\label{fol:leaf-constancy}
 h=f_0e^{-g}\text{ depends only on }z.
\end{equation}

Equations~\eqref{fol:split-weight} and \eqref{fol:leaf-constancy} imply,
on these charts,
\[
 \dd\Theta_T=0,\qquad \ddc(\phi\Theta_T)=0.
\]
For the second identity, the coefficient depending only on $z$ causes
no derivative after wedging with $\dd z\wedge\dd\bar z$, while the
remaining $g+\bar g$ is pluriharmonic in the leaf directions.
The first- and second-derivative versions of
\eqref{fol:cutoff-estimate} extend both identities across $A$.

Finally choose a smooth reference metric on $L$, with local weights
$\psi$ and curvature $\theta$.  The difference $u=\phi-\psi$ is a
global function downstairs.  The local pushforwards defining
$u\Theta_T$ agree: upstairs they are locally integrable by the
exponential estimates, they agree on the good open, and they have
no analytic-subset mass.  Closedness and the last displayed identities
give the global current identity
\begin{equation}\label{fol:potential-identity}
 \ddc(u\Theta_T)=-\theta\wedge\Theta_T.
\end{equation}
It proves \eqref{fol:Theta-orthogonality}.
\end{proof}

\subsection{Hodge index and the space of positive currents}

\begin{lemma}\label{fol:nonnegative-square}
Let $X$ be a smooth compact K\"ahler fourfold, with K\"ahler form
$\omega$, and let $R$ be a positive closed $(1,1)$-current with no
divisorial mass.  Then
\[
 \int_X [R]^2[\omega]^2\geq0.
\]
\end{lemma}

\begin{proof}
Regularization with analytic singularities gives currents
$R_k\in[R]$ and numbers $\epsilon_k\downarrow0$ such that
$R_k\geq-\epsilon_k\omega$ and
$\nu(R_k,x)\leq\nu(R,x)$ at every point; see
\cite[Theorem~2.1(ii)]{BoucksomDivisorialZariski}.
In particular the analytic pole sets of $R_k$ have codimension at least
two.  Resolve these pole sets by a projective modification
$f_k\colon X_k\to X$.  For
$\beta_k=[R]+\epsilon_k[\omega]$, the pulled positive current has a
decomposition
\[
 f_k^*\beta_k=P_k+[D_k],
\]
where $D_k$ is an effective exceptional real divisor and $P_k$ is nef.
With the usual smooth-remainder version of analytic regularization,
$P_k$ is represented by a smooth semipositive form.  The
bounded-remainder version gives a positive residual current with bounded
potentials, which is nef by a further regularization.

Since $f_{k*}[D_k]=0$, the projection formula gives
\[
 \int_{X_k} f_k^*\beta_k\,[D_k]\,(f_k^*[\omega])^2=0.
\]
Consequently
\begin{equation}\label{fol:square-resolution}
 \int_X\beta_k^2[\omega]^2
 =
 \int_{X_k} f_k^*\beta_k\,P_k\,(f_k^*[\omega])^2
 \geq0.
\end{equation}
The inequality pairs a pseudo-effective class with three nef classes.
It follows by approximating the nef factors by K\"ahler classes and
pairing with a positive current.  Passing to the limit proves the lemma.
In particular, no positivity of the self-intersection of the
exceptional divisor has been used.
\end{proof}

For a K\"ahler form $\omega$ on $M$, write
\[
 Q_\omega(\xi,\zeta)=\int_M\xi\zeta[\omega]^2
 \qquad(\xi,\zeta\in H^{1,1}(M,\R)).
\]
The Hodge index theorem says that this quadratic form has signature
$(1,h^{1,1}-1)$.  We use its following elementary consequence:
two nonzero classes with nonnegative square, positive pairing with
$[\omega]$, and zero mutual pairing must both lie on the same isotropic
ray.

Choose the zero-Lelong current supplied by
Lemma~\ref{met:zero-lelong}, and let $\Theta$ be the current of
Lemma~\ref{fol:transverse-current}.  The nef class $\alpha$ has
nonnegative square, and Lemma~\ref{fol:nonnegative-square} applies to
$\Theta$.  Both classes have positive mass against $\omega^3$;
\eqref{fol:nonzero-alpha} is used here.  Their mutual pairing is zero by
\eqref{fol:Theta-orthogonality}.  Thus
\begin{equation}\label{fol:isotropic-ray}
 Q_\omega(\alpha,\alpha)=0,\qquad [\Theta]=c\alpha
 \quad\text{for some }c>0.
\end{equation}
In fact,
\begin{equation}\label{fol:alpha-square-zero}
 \alpha^2=0\quad\text{in }H^4(M,\R).
\end{equation}
To check the stronger assertion, choose K\"ahler representatives of
$\alpha+\epsilon[\omega]$.  Their positive squares have mass against
$\omega^2$ tending to zero by \eqref{fol:isotropic-ray}.  Their
cohomology classes tend to $\alpha^2$, whereas their currents tend
weakly to zero, because their positive masses tend to zero.  Thus the
limiting cohomology class is zero.

\begin{lemma}\label{fol:all-zero-Lelong}
Every positive closed current in $\alpha$ has zero Lelong numbers
everywhere.
\end{lemma}

\begin{proof}
Suppose $T\in\alpha$ has a positive Lelong number at a point, and let
$f\colon X\to M$ be its blowup.  The pullback has a positive generic
Lelong number on the exceptional divisor.  Write its Siu decomposition
as
\begin{equation}\label{fol:Siu-decomposition}
 f^*T=R+\sum_j a_j[D_j],
 \qquad a_j>0,
\end{equation}
where $R$ has no divisorial mass and the sum is not zero.
Fix a K\"ahler form $\omega_X$ on $X$, and set $\alpha_X=f^*\alpha$.
The class $\alpha_X$ is nef, nonzero, and
$\alpha_X^2=0$.  Pairing \eqref{fol:Siu-decomposition} with
$\alpha_X[\omega_X]^2$ gives zero.  All terms are nonnegative,
since a nef class paired with a positive current and two K\"ahler
classes has nonnegative intersection.  Therefore
\[
 Q_{\omega_X}(\alpha_X,[R])=0,\qquad
 Q_{\omega_X}(\alpha_X,[D_j])=0
 \quad\text{for every }j.
\]
Lemma~\ref{fol:nonnegative-square} and Hodge index give
$[R]=b\alpha_X$ with $b\geq0$; this includes $b=0$ when $R=0$.
Comparing masses in \eqref{fol:Siu-decomposition} gives $b<1$,
because the divisorial sum has positive mass.  Hence
\begin{equation}\label{fol:divisor-span}
 (1-b)\alpha_X=\sum_j a_j\,c_1(\OO_X(D_j)).
\end{equation}

The real span of the integral classes $c_1(\OO_X(D_j))$ is a
finite-dimensional, hence closed, subspace.  The convergent series in
\eqref{fol:divisor-span} therefore places the rational class
$c_1(f^*L)=\alpha_X/(2\pi)$ in that real span.  Choose a finite subset
of these integral classes forming a basis of the span.  Solving the
corresponding rational linear system shows that a rational vector in
their real span is in their rational span.  Thus, for some rational
divisor $D$ with finite support,
\[
 c_1(f^*L)=c_1(\OO_X(D))
 \quad\text{in }H^2(X,\Q).
\]
After clearing denominators and integral torsion, the difference line
bundle has zero integral first Chern class.  By
Lemma~\ref{fol:virtual-vanishing}, $q(X)=0$, so the exponential sequence
makes this difference line bundle trivial.  A positive multiple of
$f^*L$ consequently has a nonzero meromorphic section.  The discrepancy
identity for $X\to M\to Y$ converts it to a meromorphic section of a
positive multiple of $K_X$, contradicting
Lemma~\ref{fol:virtual-vanishing}.
\end{proof}

The normalized-current fixed-point construction is inspired by Touzet's
method \cite{Touzet2014}.  The singular extension, including continuity
on the higher charts, is proved here.

\begin{lemma}\label{fol:fixed-point}
There is a positive current $T\in\alpha$ such that, on all the smooth
index charts constructed in Lemma~\ref{fol:transverse-current},
\begin{equation}\label{fol:fixed-point-equation}
 \ii\partial\bar\partial\phi
 =\ii b e^{-\phi}\eta\wedge\bar\eta
\end{equation}
for one constant $b>0$.  On these charts $\phi$ is smooth and
\begin{equation}\label{fol:eta-derivative}
 \dd\eta=\partial\phi\wedge\eta.
\end{equation}
\end{lemma}

\begin{proof}
Fix a smooth reference weight for $L$, with curvature $\theta$, and a
smooth probability volume on $M$.  The set
\[
 \mathcal C_\alpha=
 \{T\geq0:\ T\text{ closed of type }(1,1),\ [T]=\alpha\}
\]
is nonempty, compact, and convex in the weak topology of currents.
Its mass against $\omega^3$ is fixed.  Write
$T=\theta+\ddc u_T$, uniquely normalized by
$\int_Mu_T\,\dd V=0$.  Thus all local weights used for $T$ have
a fixed common additive normalization.

By Lemma~\ref{fol:all-zero-Lelong}, the construction of
Lemma~\ref{fol:transverse-current} applies to every $T\in\mathcal C_\alpha$.
Applying Hodge index as in \eqref{fol:isotropic-ray} shows that
$[\Theta_T]$ is a positive multiple of $\alpha$.  Define
\begin{equation}\label{fol:fixed-point-map}
 \mathcal F(T)=
 \frac{\int_M\alpha[\omega]^3}{\int_M\Theta_T\wedge\omega^3}\,\Theta_T.
\end{equation}
This is a selfmap of $\mathcal C_\alpha$.

We verify continuity, including on the higher charts.  If $T_j\to T$
weakly, compactness for quasi-plurisubharmonic functions with fixed lower
curvature bound and the chosen normalization gives
$u_{T_j}\to u_T$ in $L^1$.  Indeed every subsequential limit has the same
$\ddc$ and normalization as $u_T$.  On a coordinate chart the metric
weights are plurisubharmonic and converge in $L^1$.
All complex singularity exponents of the limit are infinite, by
Lemma~\ref{fol:all-zero-Lelong} and Skoda integrability.
The effective semicontinuity theorem
\cite[Theorem~0.2(2)]{DemaillyKollar} consequently gives locally uniform
bounds for $e^{-s\phi_j}$ for every fixed $s>0$, and in fact convergence
of these exponentials in $L^1$ on smaller neighborhoods.

Apply \eqref{fol:Jacobian-Holder} on the fixed proper chart diagrams.
It gives locally uniform upstairs bounds for every inverse exponential
power as well.  After extraction, the weights converge almost everywhere
off the exceptional and critical analytic sets.  Pullback preserves
almost-everywhere convergence there, because the chart maps are local
biholomorphisms; the omitted analytic sets have measure zero.
Uniform $L^r$ bounds for some $r>1$ then give convergence in $L^1$
of the densities defining $\Theta_{T_j}$ upstairs.  Proper
pushforward gives $\Theta_{T_j}\to\Theta_T$ downstairs.  The argument
applies to each subsequence and so proves continuity for the whole
sequence.  Their masses also converge, and the limiting mass is positive.
Thus $\mathcal F$ is continuous.  The Schauder--Tychonoff fixed-point
theorem yields a fixed point.

Let $b$ be the positive normalization factor in
\eqref{fol:fixed-point-map} at this fixed point.  Its equality
$T=b\Theta_T$ holds upstairs on the good locus, giving
\eqref{fol:fixed-point-equation}.  Both sides extend with no
analytic-subset mass, so the equality holds on each entire smooth
higher chart.

Taking a smooth Euclidean trace on a smaller coordinate ball gives a
scalar Poisson equation whose right side belongs to every finite $L^r$,
because $\eta$ is holomorphic and all inverse exponential moments of
$\phi$ are finite.  Interior elliptic regularity gives
$\phi\in W^{2,r}_{\mathrm{loc}}$ for every finite $r$.
Choosing $r>8$ gives continuous first derivatives, after which the
semilinear equation and ordinary elliptic bootstrap make $\phi$ smooth.
On a regular foliated chart,
\eqref{fol:split-weight} and \eqref{fol:leaf-constancy} give
$\eta=e^g h(z)\,\dd z$.  Therefore
$\dd\eta=\dd g\wedge\eta=\partial\phi\wedge\eta$ there.
Both sides are now smooth on the entire higher chart, and the identity
extends by density.  This proves \eqref{fol:eta-derivative}.
\end{proof}

\subsection{Spherical developing maps and boundary meridians}

Fix the current given by Lemma~\ref{fol:fixed-point}.  On a smooth higher
chart put
\begin{equation}\label{fol:unitary-forms}
 \gamma=\sqrt{b/2}\,e^{-\phi/2}\eta,\qquad
 a=\frac{\partial\phi-\bar\partial\phi}{2}.
\end{equation}
Equations~\eqref{fol:fixed-point-equation} and
\eqref{fol:eta-derivative} imply
\begin{equation}\label{fol:structure-equations}
 \dd\gamma=a\wedge\gamma,\qquad
 \dd a=-2\gamma\wedge\bar\gamma.
\end{equation}
Indeed differentiating $e^{-\phi/2}\eta$ gives the first identity, while
$\dd a=-\partial\bar\partial\phi$ gives the second.  Consequently the
matrix-valued one-form
\begin{equation}\label{fol:flat-matrix}
 A=
 \begin{pmatrix}
  a/2 & \bar\gamma\\
  -\gamma & -a/2
 \end{pmatrix}
\end{equation}
is skew-Hermitian, trace-free, and satisfies
$\dd A+A\wedge A=0$.  On a simply connected small ball it therefore
has the form $U^{-1}\dd U$ with $U$ valued in $\mathrm{SU}(2)$.

The map
\[
 F\colon x\longmapsto U(x)\C e_1\in\PP^1
\]
is holomorphic.  In fact,
$\bar\partial(Ue_1)=\tfrac12a^{0,1}Ue_1$, because $\gamma$ has type
$(1,0)$; hence the line spanned by $Ue_1$ is a holomorphic line.
This argument remains valid at zeros of $\eta$.
For the Fubini--Study normalization
$\omega_{\mathrm{FS}}=\ii\partial\bar\partial\log(1+|z|^2)$,
the usual unitary-frame computation gives
$F^*\omega_{\mathrm{FS}}=\ii\gamma\wedge\bar\gamma$.
Thus, with $\omega_{\mathrm{sph}}=2\omega_{\mathrm{FS}}$,
\begin{equation}\label{fol:spherical-pullback}
 F^*\omega_{\mathrm{sph}}=T
\end{equation}
on the higher chart, where the right side is pulled back from $M$.

Choose a connected dense open subset $S\subset M$ with analytic
complement such that $p$ identifies $S$ with a smooth open subset of
$Y$, the relevant forms have their generic ranks, and
$\boldsymbol\eta$ is nowhere zero on $S$.  The index charts are
unramified over $S$, and their resolutions can be chosen isomorphic
there.  Equation~\eqref{fol:spherical-pullback} gives local holomorphic
submersions from $S$ to $\PP^1$.  Two such submersions with the same
pullback form differ locally by a constant element of
$\mathrm{PU}(2)$.  Indeed their kernels agree, so one factors locally
through the other on a transverse disc; the resulting local holomorphic
isometry of the round sphere is the restriction of a projective unitary
transformation.  The transformation is unique on each connected
overlap because the image of a submersion contains an open subset.

These local maps and their constant transitions give a developing map
and a monodromy homomorphism
\begin{equation}\label{fol:developing-map}
 \operatorname{dev}\colon\widetilde S\to\PP^1,\qquad
 \rho\colon\pi_1(S)\to\mathrm{PU}(2).
\end{equation}
The developing map is nonconstant and is equivariant for $\rho$.

\begin{lemma}\label{fol:finite-meridians}
Let $X$ be a smooth space with a proper generically finite map to $M$
that is unramified over $S$, and let $S_X$ be the inverse image of $S$.
After any further modification supported outside $S_X$, the monodromy
of a small meridian about every prime divisor in the complement of
$S_X$ has finite order.
\end{lemma}

\begin{proof}
The assertion is local at a general point of a boundary prime $D$.
Choose a small disc $\Delta$ transverse to $D$ at a smooth point not
on another boundary component, with $\Delta^*=\Delta\setminus\{0\}$
contained in $S_X$.  Choose an original proper higher chart over a
neighborhood of the image point in $M$, take its fiber product with
$X$, and resolve a main component.  The resulting map
$g\colon W\to X$ is proper and generically finite, and is unramified
over $S_X$.  Near every point of $W$ there is an ambient holomorphic
map to $\PP^1$: compose its map to the original higher chart with the
map constructed in \eqref{fol:flat-matrix}.  Over the good open these
maps agree with the developing germs up to constant projective
unitary transformations.

A component of the inverse image of $\Delta$ dominating $\Delta$
is a curve, finite over $\Delta^*$.  Its normalization has a point over
$0$, and its local map to $\Delta$ can be written, after changing a
parameter, as $t\mapsto t^r$ for some positive integer $r$.
For a sufficiently small circle in the $t$-disc, the lifted loop stays
inside one neighborhood on which an ambient map to $\PP^1$ is
defined.  Continuing this ambient map around the lifted loop returns
the same germ.  On the punctured part, $g$ is locally biholomorphic,
so this is the analytic continuation of the original developing germ
around the $r$-th power of the meridian.

A projective unitary transformation fixing that germ must be the
identity, since it is a submersive ambient germ on the good open.
Thus the meridian has monodromy whose $r$-th power is the identity.
It is essential here to use the ambient germ: even if the transverse
disc happens to be tangent to the foliation, the argument does not
infer identity of holonomy merely from its action on the values of
the map along that disc.  The same fiber-product construction applies
after the further modifications specified in the statement.
\end{proof}

We record the compactification fact in the precise analytic category
needed here.

\begin{lemma}\label{fol:cover-compactification}
Let $X$ be a smooth compact K\"ahler manifold and $D$ a simple normal
crossing divisor.  Every finite unramified cover of $X\setminus D$
extends to a finite normal cover of $X$.  It has a smooth compact
K\"ahler resolution which is unchanged over the original cover, and
whose complement of that cover can be made a simple normal crossing
divisor.
\end{lemma}

\begin{proof}
Near a point of $D$, choose a polydisc on which its complement is
$(\Delta^*)^r\times\Delta^{n-r}$.  Each connected finite cover is
described by a finite-index subgroup of its fundamental group
$\Z^r$.  Such a subgroup contains $N\Z^r$ for some $N>0$.
The cover is consequently dominated by the coordinate power cover
\[
 (t_1,\ldots,t_r,w)\longmapsto(t_1^N,\ldots,t_r^N,w).
\]
It extends across the coordinate hyperplanes as the normal quotient
of the full polydisc power cover by the corresponding finite subgroup
of its deck group.  The normal finite extensions glue uniquely:
an isomorphism on the dense punctured locus extends between the
normal finite algebras, equivalently by their integral closures.
This gives a finite normal cover $\overline X\to X$.

For completeness, a finite source over a compact K\"ahler space is
K\"ahler in the local-embedding sense.  Over finitely many small base
neighborhoods choose relative embeddings of the finite source into
products with affine spaces, with relative coordinates $w_a$.
Choose a smooth partition of unity $\lambda_a$ on the base and put
\[
 h=\sum_a(\lambda_a\circ f)|w_a|^2
\]
on the source, extending each summand by zero outside its support.
This is a smooth function on the complex space.  At a point where
$\lambda_a>0$, use the corresponding relative embedding.
The other relative coordinate functions have local holomorphic
extensions to its ambient space.  On vertical tangent directions the
Levi form of $h$ is
$\sum_a\lambda_a|\dd w_a|^2$ and is strictly positive.  The terms
arising from derivatives of the partition have a base-direction
factor.  A sufficiently large multiple of the pulled base K\"ahler
form dominates their horizontal and mixed contributions.  Compactness
allows one such multiple on the finite cover.  Thus
\[
 C f^*\omega_X+\ddc h
\]
has strictly plurisubharmonic local potentials in the chosen ambient
embeddings and defines a K\"ahler form on the normal source.
Finally take a projective resolution and then an embedded resolution
of the boundary, both unchanged over the smooth covered open.
Projective modifications of compact K\"ahler spaces remain K\"ahler,
by the usual relative ample curvature construction.  This proves the
claim.
\end{proof}

Resolve $M\setminus S$ by a projective modification, unchanged on $S$,
so that the complement is a simple normal crossing divisor.
The fundamental group of $S$ is finitely generated: a complement of
this kind in a compact smooth manifold has the homotopy type of a
finite complex, or one may retract onto a compact manifold with
corners obtained by removing sufficiently small tubular neighborhoods.

The group $\mathrm{PU}(2)$ is linear, for example through the adjoint
embedding into $\mathrm{GL}_3(\C)$.  Selberg's lemma therefore gives
a finite-index torsion-free subgroup of $\rho(\pi_1(S))$.
Take the connected finite unramified cover $S_0\to S$ associated with
its inverse image in $\pi_1(S)$.  By
Lemma~\ref{fol:cover-compactification}, it is contained in a smooth
compact K\"ahler manifold $M_0$, with simple normal crossing
complement, and there is a proper generically finite morphism
$M_0\to M$.  Its monodromy image is torsion-free.

Every boundary meridian of $M_0\setminus S_0$ has finite-order
monodromy by Lemma~\ref{fol:finite-meridians}, hence trivial
monodromy.  The inclusion $S_0\hookrightarrow M_0$ induces a
surjection of fundamental groups whose kernel is normally generated
by those meridians.  One way to see both assertions is to put loops
and homotopies in general position relative to the normal crossing
divisor: loops avoid it, and a homotopy meets its smooth part in
finitely many transverse points, each contributing a meridian; it
avoids the intersections of components.  Consequently the
representation descends to
\begin{equation}\label{fol:compact-representation}
 \rho_0\colon\pi_1(M_0)\to\mathrm{PU}(2).
\end{equation}

\subsection{Finite holonomy and the contradiction}

\begin{lemma}\label{fol:finite-unitary-image}
Let $X$ be a connected smooth compact K\"ahler manifold such that
$a(X)=0$ and $q(X')=0$ for every connected finite \'etale cover
$X'\to X$.  Every representation
$\pi_1(X)\to\mathrm{PU}(2)$ with torsion-free image has trivial image.
\end{lemma}

\begin{proof}
Let $\Gamma$ be the image, and take its complex Zariski closure in
$\mathrm{PSL}_2(\C)$, regarded as a linear algebraic group through
the adjoint representation.

If this closure is all of $\mathrm{PSL}_2(\C)$, apply the
semisimple Shafarevich theorem
\cite[Theorem~1]{CCE}.  Its hypotheses are exactly a smooth compact
K\"ahler source and a representation Zariski dense in a semisimple
linear group; for torsion-free image the Shafarevich base is normal
projective of general type, and the representation factors through
a smooth model of that base.  Algebraic dimension zero forces this
base to be a point: a positive-dimensional projective base would
supply a nonconstant meromorphic function on $X$.  Factorization
through a point makes $\Gamma$ trivial, contradicting its supposed
Zariski density.

If the Zariski closure is proper, its identity component is solvable,
by the classification of proper connected algebraic subgroups of
$\mathrm{PSL}_2(\C)$.  Thus $\Gamma$ is virtually solvable.  Pass to
a finite-index subgroup whose image $\Gamma_0$ is solvable and let
$X_0\to X$ be the corresponding finite \'etale cover.
For every further finite-index subgroup $H$ of $\pi_1(X_0)$, its
abelianization is finite: it is the first integral homology of a
compact K\"ahler finite cover with $q=0$, hence is finitely generated
of rank $2q=0$.

Now induct along the derived series of $\Gamma_0$.
The quotient $\Gamma_0/\Gamma_0'$ is an image of
$\pi_1(X_0)^{\mathrm{ab}}$, so it is finite.
Its preimage has finite index in $\pi_1(X_0)$; that subgroup again
has finite abelianization, making $\Gamma_0'/\Gamma_0''$ finite.
Repeating gives finite successive indices down to the identity,
because the derived series has finite length.
Thus $\Gamma_0$, and then $\Gamma$, is finite.
A finite torsion-free group is trivial.
\end{proof}

Lemma~\ref{fol:virtual-vanishing} supplies the hypotheses of
Lemma~\ref{fol:finite-unitary-image} for $M_0$ and all its finite
\'etale covers.  The representation
\eqref{fol:compact-representation} is therefore trivial.  The
developing map descends to a single-valued nonconstant holomorphic map
\[
 F_0\colon S_0\longrightarrow\PP^1.
\]

We finish by proving its meromorphic extension, rather than appealing
to compactness of the target.  Over a neighborhood in $M_0$, form the
proper generically finite higher charts used in
Lemma~\ref{fol:finite-meridians}.  Each small ball upstairs has an
ambient sphere map obtained by composition from
\eqref{fol:flat-matrix}.  On the good part of the ball this map and
the lift of $F_0$ have the same transverse round metric, so they differ
by a constant projective unitary transformation.  The good part,
being the complement of a proper analytic subset in a ball, is
connected; uniqueness on submersive germs makes that transformation
constant throughout it.  After this transformation the ambient map
extends the lift of $F_0$.  The extensions agree on intersections
by density and glue to a holomorphic map on the higher chart.

Choose a scalar projective coordinate on $\PP^1$.  Composing with
these extensions gives a meromorphic function on each proper
generically finite chart.  Its meromorphic trace, divided by the
chart degree, is a meromorphic function on the base neighborhood.
To define the trace, first descend through the proper modification
in the chart's Stein factorization and then take the ordinary trace
of the finite normal map.  On the good open every branch is the
same pullback of the chosen coordinate of $F_0$, so its trace is
exactly the degree times that coordinate.  The resulting local
meromorphic functions on $M_0$ consequently agree on overlaps and
give a global meromorphic extension of the coordinate of $F_0$.
It is nonconstant, contradicting $a(M_0)=0$ from
Lemma~\ref{fol:virtual-vanishing}.

This contradiction proves Proposition~\ref{fol:empty-case}.
